%TCIDATA{LaTeXparent=0,0,EOMN.tex}
                      

\section{Commutative Group ${\frak G}_{\ker }$\protect\bigskip }

We can look for a unitary subgroup of ${\cal U}({\cal H}^{N})$ such that it
maps equivalence classes $\left[ f\right] _{N}$ to themselves and sends
positive operators to positive operators. If it exists it has the properties%
\begin{eqnarray}
&&Tr\{\Phi ^{\dagger }\Phi UYU^{\dagger }\}=Tr\{\Phi ^{\dagger }\Phi
Y\}\quad \forall \,Y\in {\cal B}_{1}({\cal H}^{N})  \nonumber \\
&\Rightarrow &Tr\{\Phi ^{\dagger }\Phi UD^{N}U^{\dagger }\}=Tr\{\Phi
^{\dagger }\Phi D^{N}\}\quad \forall \,D^{N}\in {\cal B}_{1}({\cal H}^{N}) 
\nonumber \\
&\Rightarrow &\left\langle U\Psi |\Phi ^{\dagger }\Phi U\Psi \right\rangle
=\left\langle \Psi |\Phi ^{\dagger }\Phi \Psi \right\rangle \quad \forall
\Psi \in {\cal H}^{N}  \nonumber \\
&\Leftrightarrow &U^{\dagger }\Phi ^{\dagger }\Phi U=\Phi ^{\dagger }\Phi
\Longleftrightarrow \left[ \Phi ^{\dagger }\Phi ,U\right] =0  \nonumber \\
&\Longleftrightarrow &U=f\left( \Phi ^{\dagger }\Phi \right) \Rightarrow
U=\exp i\int \omega \left( \underline{\kappa }\right) \prod_{1\leq \iota
\leq N}\Phi \left( \kappa _{i}\right) ^{\dagger }\Phi \left( \kappa
_{i}\right) d\underline{\kappa }
\end{eqnarray}%
i.e. it is a unitary commutative subgroup, ${\frak G}_{\Xi }$, of ${\cal U}(%
{\cal H}^{N})$ and it is diagonal in the coordinate representation, {\bf %
which is its defining representation}. Hence, trivially, it has the
representation, ${\cal R}$, in the space of bounded operators acting in $%
{\cal H}^{N}$%
\begin{eqnarray}
{\cal R} &:&{\cal G}_{\Xi }\longrightarrow {\cal B}\left( {\cal H}^{N}\right)
\nonumber \\
{\cal R}\left( g\right) &=&\exp i\int \omega _{g}\left( \underline{\kappa }%
\right) \prod \Phi \left( \kappa _{i}\right) ^{\dagger }\Phi \left( \kappa
_{i}\right) d\underline{\kappa };\quad g\in {\cal G}_{\Xi }  \label{unitary}
\end{eqnarray}%
This representation has the property that ${\cal R}\left( {\cal G}_{\Xi
}\right) :$ $\left[ f\right] _{N}\longrightarrow \left[ f\right] _{N}\quad
\forall f$ \ i.e. every ${\cal R}\left( g\right) $ maps all the members of
an equivalence to the same equivalence class and the same representation
works for all equivalence classes, explicitly $UYU^{\dagger }\in \left[ f%
\right] _{N}\Longleftrightarrow Y\in \left[ f\right] _{N}$ for $U\in {\cal R}%
\left( {\cal G}_{\Xi }\right) $. (Note a pure state $\left| \Psi
\right\rangle $ defines an element of an equivalence class by $\left| \Psi
\right\rangle \left\langle \Psi \right| $).

${\cal R}\left( {\cal G}_{\Xi }\right) $ does not produce all pure states
with the same density as a reference pure state, when acting on that
reference state. This is proved by noting that $\langle $\underline{$\kappa $%
}${\bf |}{\cal U}\left( \omega \right) \Psi \rangle =e^{i\omega \left( 
\underline{\kappa }\right) }\Psi \left( \underline{\kappa }\right) $ \ does
not range over all pure state wavefunctions that have the same density for a
fixed $\Psi $ and varying $\omega .$ This can be seen from the observations
that $e^{{\cal V}\left( \underline{\kappa }\right) }\Psi \left( \underline{%
\kappa }\right) $ has the same density as $\Psi \left( \underline{\kappa }%
\right) $ when $e^{{\cal V}\left( \underline{\kappa }\right) }=\left( 1+\phi
\left( \underline{\kappa }\right) \right) ^{\frac{1}{2}}$ and $\phi \left( 
\underline{\kappa }\right) $ is strongly orthogonal to $\left| \Psi \left( 
\underline{\kappa }\right) \right| ^{2}$ and noting that $e^{{\cal V}\left( 
\underline{\kappa }\right) }\Psi \left( \underline{\kappa }\right) $ cannot
be obtained from $\Psi \left( \underline{\kappa }\right) $ by multiplying by 
$e^{i\omega \left( \underline{\kappa }\right) }.$ Thus the action of the
unitary commutative subgroup does not produce all pure states with the same
density.

If one drops the unitary requirement {\it it is} possible find
transformations, {\it that are diagonal in the coordinate representation, }%
which produce all pure states with a given density, $\gamma $, from a pure
reference state $\left| \Psi \right\rangle $ of that density{\it . }These
transformations are of the form%
\begin{equation}
\breve{T}=\exp \left\{ \int \left\{ {\cal V}\left( \underline{\kappa }%
\right) +i\omega \left( \underline{\kappa }\right) \right\} \prod_{1\leq
\iota \leq N}\Phi \left( \kappa _{i}\right) ^{\dagger }\Phi \left( \kappa
_{i}\right) d\underline{\kappa }\right\} {\bf \quad \Rightarrow }\breve{T}%
\in GL\left( {\cal H}^{N}\right)
\end{equation}%
where $e^{{\cal V}\left( \underline{\kappa }\right) }=\left( 1+\phi \left( 
\underline{\kappa }\right) \right) ^{\frac{1}{2}}$ and $\phi \left( 
\underline{\kappa }\right) $ is strongly orthogonal to $\left| \Psi \left( 
\underline{\kappa }\right) \right| ^{2}$. The proof of this assertion is by
construction and by noting that these transformations produce all pure
states from any reference pure state. This can be seen by noting that

\begin{equation}
\left[ \breve{T}\Psi \right] \left( \underline{\kappa }\right) =e^{{\cal V}%
\left( \underline{\kappa }\right) +i\omega \left( \underline{\kappa }\right)
}\Psi \left( \underline{\kappa }\right)
\end{equation}%
and by considering the local norm/phase form of the initial and final
wavefunctions%
\begin{eqnarray}
\langle \underline{\kappa }\left| \Psi \right\rangle &=&N\left( \underline{%
\kappa }\right) e^{i\theta \left( \underline{\kappa }\right) } \\
\langle \underline{\kappa }\left| \breve{T}\Psi \right\rangle &=&N^{\prime
}\left( \underline{\kappa }\right) e^{i\theta ^{\prime }\left( \underline{%
\kappa }\right) }
\end{eqnarray}%
These transformations form{\it \ }{\bf a defining representation of a
commutative subgroup}{\it \ }${\frak G}_{fib}={\cal G}_{\ker }\times {\cal G}%
_{\Xi }$. This is the{\bf \ }abstract group whose representation ${\cal R}%
\left( {\frak G}_{fib}\right) $ maps each equivalence class 
\begin{equation}
\lbrack f]_{N}\equiv \left\{ \left( \breve{k},\breve{\Gamma}\left( f\right)
\right) ;\breve{k}\in \ker \Xi ,\text{ }\breve{\Gamma}\left( f\right) \in
\ker \Xi ^{c},\text{fixed }f\right\}
\end{equation}%
to itself for {\bf any }$f$ i.e. 
\begin{equation}
{\cal R}\left( g\right) \left( \breve{k},\breve{\Gamma}\left( f\right)
\right) =\left( \breve{k}^{\prime },\breve{\Gamma}\left( f\right) \right)
;\quad g\in {\frak G}_{fib};\,{\cal R}\left( g\right) \breve{k}=\breve{k}%
^{\prime };{\cal R}\left( g\right) \breve{\Gamma}\left( f\right) =\breve{%
\Gamma}\left( f\right) \,
\end{equation}%
The representation is 
\begin{eqnarray}
{\cal R} &:&{\frak G}_{fib}\longrightarrow {\cal B}\left( {\cal H}^{N}\right)
\nonumber \\
{\cal R}\left( g\right) &=&\exp \int \left\{ {\cal V}\left( \underline{%
\kappa }\right) +i\omega \left( \underline{\kappa }\right) \right\} \prod
\Phi \left( \kappa _{i}\right) ^{\dagger }\Phi \left( \kappa _{i}\right) d%
\underline{\kappa ;}\qquad g\in {\frak G}_{fib};
\end{eqnarray}%
where $\left\{ {\cal V}\left( \underline{\kappa }\right) \right\} $ have the
property 
\begin{equation}
{\cal V}\left( \underline{\kappa }\right) \Gamma \left( f\right) \left( 
\underline{\kappa }\right) =0\quad a.e.
\end{equation}%
wrt $\Gamma \left( f\right) \left( \underline{\kappa }\right) \quad \forall
f\in L_{1}\left( Z\right) .$ Explicitly $TYT^{\dagger }\in \left[ f\right]
_{N}\Longleftrightarrow Y\in \left[ f\right] _{N}$ for $T\in {\cal R}\left( 
{\frak G}_{fib}\right) $. The group ${\frak G}_{fib}$ does {\it not }act
homogeneously in $\left[ f\right] _{N}$, i.e. it does not generate all of $%
\left[ f\right] _{N}$ from a reference element, but it does generate all
pure states in $\left[ f\right] _{N}$ from a reference pure state in $\left[
f\right] _{N}$ i.e. all pure states lie on a single orbit of ${\frak G}%
_{fib} $ in $\left[ f\right] _{N}$ in contrast to orbits of ${\cal G}_{\Xi }$%
. ({\bf Note }that{\bf \ }the groups ${\cal G}_{\Xi }$ and ${\frak G}_{fib}$
act in ${\frak B}_{vect}$ through their representations as operators in $%
{\cal B}_{1}\left( {\cal H}^{N}\right) $.) Agian {\bf one} representation
suffices for {\bf all }equivalence classes.

The set of pure state wavefunctions $\left\{ \Phi _{\gamma }\left( 
\underline{\kappa }\right) \right\} $ associated with pure states in an
equivalence class $\left[ \gamma \right] _{N}$ can all be expressed in terms
of a reference wavefunction $\Psi \left( \underline{\kappa }\right) $, and
the non redundant parameterizing functions ${\cal V}$ and $\omega $ as $%
\left\{ \Phi \left( \underline{\kappa }\right) =e^{{\cal V}\left( \underline{%
\kappa }\right) +i\omega \left( \underline{\kappa }\right) }\Psi \left( 
\underline{\kappa }\right) =\left\langle \underline{\kappa }\left| e^{{\cal V%
}+i\breve{\omega}}\right. \Psi \right\rangle \right\} .$ If the reference
state is normalized then the transformed states $\left| e^{{\cal V}+i\breve{%
\omega}}\Psi \right\rangle $ $\equiv \left| \Psi \left( {\cal V},\omega
\right) \right\rangle $ are normalized. The question of {\bf whether there
exists} a set of unitary operators $\subset U\left( {\cal H}^{N}\right)
\subset {\cal B}\left( {\cal H}^{N}\right) $ that form a group and map $%
\left\{ \left| \Psi \left( {\cal V},\omega \right) \right\rangle
\left\langle \Psi \left( {\cal V},\omega \right) \right| \right\} $ $\subset %
\left[ \gamma \right] _{N}$ to themselves is an open question. Clearly such
a set of operators would have be noncommutative and the set would depend on $%
\gamma $ and correspond to different (but probably equivalent) unitary
representations of some abstract group.

Luckily we do not need a unitary solution so the commutative non unitary one
obtained above suffices and has the added nice property of mapping diagonal
operators in the coordinate representation to diagonal operators and
nondiagonal ones to nondiagonal ones, as well as being commutative.
